\begin{abstract}
Recent advances in brain decoding work has shown the potential of vision transformers (ViTs) to learn mappings from neural representation to visual features. Here we present CortexDecoder, a brain decoding framework mapping fMRI flat-maps to images viewed by subjects. Flat-maps are neural responses projected onto cortical surfaces and then unfolded into 2D images. We demonstrate that (1) flat-maps preserve sufficient signal to achieve high performance in brain decoding; raw fMRI data has noisy signals and high-dimensionality which makes decoding a difficult task. (2) Reshaping raw fMRI data into flat-maps allows for clean patch-based tokenization. This allows us to harness state-of-the-art Vision Transformers to learn meaningful neural representations and preserve spatial encoding. To do this, we pass flat-map patches into a pretrained CortexMAE encoder to obtain brain tokens. A lightweight cross-attention bridge then learns to map these brain tokens to image features. Our cross-attention bridge is split into two modules; the semantic module for predicting CLIP embeddings and the lower-level module to generate low-level image reconstruction. Finally, a fine-tuned SDXL (Stable Diffusion XL) takes this low-level reconstruction image as the initial image. Conditioned by the CLIP semantic embedding, the SDXL produces the final reconstruction image. In single-subject evaluations, CortexDecoder achieves performance comparable to state-of-the-art brain decoding methods, demonstrating the efficacy of flat-map based neural representations for visual reconstruction.
\end{abstract}
